\hspace{3mm}

\centerline{\bf \Huge Домашнее задание.}
\centerline{\bf \Huge Геометрические прогрессии}

\begin{flushright}
    \scriptsize{Даже самый быстрый глаз можно обмануть. \\Гото. Хантер}
\end{flushright}    


\problem{1!} Найдите сумму первых десяти членов геометрической прогрессии:
$1 - \dfrac{3}{2} + \dfrac{9}{4} - \dfrac{27}{8} + \dotsc$

\hspace{3mm}

\problem{2!} Найдите сумму бесконечно убывающей геометрической прогрессии: $1 + \dfrac{1}{7} + \dfrac{1}{49} + \dotsc$

\hspace{3mm}

\problem{3} Известно, что каждый Лосось в 13:35 съедает ровно одно киви и создает двух новых таких же Лососей. Сколько киви было съедено Лососями с 1 по 30 марта, если 5 марта в 12:00 в пруду плавали 27 Лососей?

\hspace{3mm}

\problem{4} Киви НеГерман решил пешком дойти до Москвы из Элисты (расстояние между городами 1301 км). В первый день НеГерман прошел 650 км, во второй - 325 км, в третий - 162,5 км и так далее. Сколько потребуется дней НеГерману, чтобы дойти до Москвы?

\hspace{3mm}

\problem{5}\textbf{Бонусная задача!(Гроб от РВ)} Анджур Аркадьевич выписал на доске все числа от $n+1$ до $2n$ и перемножил их. Расул Владимирович увидев это, начал считать на сколько нулей оканчивается это произведение и получил  $\biggl[ \dfrac{11n}{36} \biggr]$(квадратные скобки означают целую часть от числа. Например $\bigl[ 0,7 \bigr] = 0$ или  $\bigl[ 22,2 \bigr] = 22$). Ошибся ли он?
